\section{Matching coefficients}
After performing the continuum limit of the hadronic matrix element
for a fixed $t>0$ value of the flow time, the connection with
physical renormalized matrix element at $t=0$ is obtained
with a SFTX~\cite{Luscher:2013vga}.
The SFTX of the twist-2 operators $O_n(x,t)$, in general contains power
divergences with terms like $1/t^{m} \delta_{\mu_i \mu_j}$,
where $m=(d_n - d_p)/2$ depends on the dimensions of $O_n$, $d_n$,
and on the dimension of the lower-dimensional operator, $d_p$.
These terms, though, are classified using the continuum O($4$) symmetry and,
as a result, the SFTX of the symmetrized and traceless operators,
$\widehat{O}_n$, receives
contributions only from the corresponding traceless operators at $t=0$.

We renormalize the flowed fermion fields adopting
ringed fields, introduced in the previous section, and define the ringed twist-2
operators
\bea
\widehat{\rO}_n^{rs}(x,t) &=&
\rchibar^r(x,t) \gmuopen1 \lrDmu2 \cdots \lrDmuclose{n} \rchi^s(x,t) - \nonumber \\
&-& \text{terms~with~}\delta_{\mu_i \mu_j}\,,
\label{eq:flowed_t2_ringed}
\eea
where the subtraction guarantees that the resulting operator has vanishing traces.
The SFTX of $\widehat{\rO}_n^{rs}$ is
\be
\widehat{\rO}_n^{rs}(t) = c_n(t,\mu)\widehat{O}_{n,\MS}^{rs}(\mu) + \cdots\,,
\ee
where we make explicit only the dependence on the flow time and
the renormalization scale. The neglected terms are higher-order
contributions of positive powers of $t$.

The matching coefficients are determined considering
off-shell amputated one-particle irreducible (1PI) Green's functions
containing the flowed operators $\widehat{\rO}_n^{rs}(t)$.
The matching equations with massless fermions
\be
\left\langle \psi^r \widehat{\rO}_n^{rs}(t) \psibar^s \right\rangle = c_n(t,\mu)
\left\langle \psi^r \widehat{O}_{n,\MS}^{rs}(t=0,\mu) \psibar^s \right\rangle
\label{eq:matching_eq}
\ee
are solved in $D = 4 - 2 \epsilon$ at 1-loop in perturbation theory with
\be
c_n(t,\mu) =  1 + \frac{\gbar^2(\mu)}{\left(4 \pi\right)^2}c_n^{(1)}(t,\mu) +
O(\gbar^4)\,.
\label{eq:cn_1l}
\ee
The matching coefficients $c_n(t,\mu)$ depend on the renormalized coupling
and are independent of the soft scales the calculation of the matching coefficients.
In the Supplemental Material~\cite{SupplMat} we describe their calculation
and, based on the observation that they satisfy a renormalization
group equation, we also provide a resummed expression at next-to-leading log (NLL).

The final result for the matching coefficients is
\be
c_n^{(1)}(t,\mu) = C_F \left[ \gamma_n \log \left(8 \pi \mu^2 t \right) + B_n\right]\,,
\label{eq:matching_1l}
\ee
where $\gamma_n = 1 + 4 \sum_{j=2}^n \frac{1}{j} - \frac{2}{n(n+1)}$, and
\bea
B_n &=& \frac{4}{n(n+1)} + 4 \frac{n-1}{n}\log 2 + \frac{2-4 n^2}{n(n+1)}\gamma_E - \nonumber \\
&-& \frac{2}{n(n+1)}\psi(n+2) + \frac{4}{n}\psi(n+1) - 4 \psi(2) - \nonumber \\
&-& 4 \sum_{j=2}^n \frac{1}{j(j-1)} \frac{1}{2^j} \phi(1/2,1,j) - \log \left(432\right)\,.
\label{eq:finite}
\eea
In this expression $\gamma_E = 0.57221\ldots$ is the Euler's constant, $\psi(z)$
is the digamma function, and $\phi(z,s,a)$ is the Lerch transcendent defined as
$\phi(z,s,a) = \sum_{k=0}^{\infty} \frac{z^k}{(k+a)^s}$.

The expression for $\gamma_n$ matches the 1-loop anomalous dimension of the twist-2
operators~\cite{Gross:1974cs} and provides a welcome check of the calculation.
For $n>2$ the finite part, $B_n$, in Eq.~\eqref{eq:finite} provides a new result.
The result with $n=2$ was already obtained
in Ref.~\cite{Makino:2014taa} in the analysis of the energy-momentum tensor,
and we reproduce their result.

The fact that the matching of the flowed matrix elements is so simple
opens the possibility of computing moments of the parton distribution functions
of any order.
A possible operational procedure is the following: After fixing $n$,
one constructs the symmetric and traceless operator $\widehat{O}_n^{rs}$.
One proceeds calculating the matrix elements between hadron states of
$\widehat{\rO}_n^{rs}(t)$
using standard techniques based on the spectral decomposition.
To apply the decomposition
in the lattice calculation of the 3-point function,
one needs to ensure that the physical distance
between the interpolators of the
hadron states and of the flowed operators is much larger than
the flow-time radius $\sqrt{8t}$. After performing the continuum limit,
the renormalized matrix element in the $\MS$ scheme is calculated
simply by
\be
A_n^{\MS}(\mu) =
c_n(t,\mu)^{-1} A_n(t) \,.
\label{eq:x_MS}
\ee
This matching is correct only if the flowed fields are defined as ringed fields.
If the flowed fields are renormalized in a different scheme, the finite
part of the matching coefficient has to be modified.

If the matching is successful, the l.h.s of Eq.~\eqref{eq:x_MS}
should be independent of the flow time.
Violations stem from higher dimensional operators as linear terms in $t$,
and from higher order terms in the perturbative
expansion of the matching coefficients.
To mitigate these possible systematics it is important to
extend to O($\gbar^4$) the calculation presented here
and conduct a thorough numerical study of the residual
flow-time dependence of the hadronic matrix element after the matching.
